% main.tex
% A Bounded, Epistemically Labeled Verification Harness for Phenomenological Quantum-Channel Models of Curved Spacetime
%
% Compile:  latexmk -pdf main.tex   (uses biber for the bibliography)
%

\documentclass[11pt,a4paper]{article}

% --- Packages ---------------------------------------------------------------
\usepackage[utf8]{inputenc}
\usepackage[T1]{fontenc}
\usepackage{lmodern}
\usepackage[margin=1in]{geometry}
\usepackage{amsmath, amssymb, amsthm}
\usepackage{microtype}
\usepackage{enumitem}
\usepackage{hyperref}
\usepackage{xcolor}
\usepackage{booktabs}
\usepackage{tabularx}
\usepackage[backend=biber,style=alphabetic,sorting=nyt]{biblatex}

\addbibresource{references.bib}

\hypersetup{
  colorlinks=true,
  linkcolor=black,
  citecolor=black,
  urlcolor=blue!50!black,
  pdftitle={A Bounded, Epistemically Labeled Verification Harness for Phenomenological Quantum-Channel Models of Curved Spacetime},
  pdfauthor={Yun Kwansub}
}

% --- Reproduction/Verification environment ---------------
\newtheoremstyle{reproduction}%
  {1.2em}{1.2em}%        % top/bottom spacing
  {\itshape}{}%          % body / indent
  {\bfseries}{.}%        % heading / punctuation
  {.5em}%                % space after head
  {}                     % default head spec
\theoremstyle{reproduction}
\newtheorem{reproduction}{Reproduction}[section]
\let\verification\reproduction
\let\endverification\endreproduction

% Macros
\DeclareUrlCommand{\tname}{\urlstyle{tt}%
  \def\UrlBreaks{\do\.\do\/\do\_\do\-\do\:\do\(\do\)}}
\newcommand{\Q}{\mathbb{Q}}
\newcommand{\Z}{\mathbb{Z}}
\newcommand{\C}{\mathbb{C}}
\newcommand{\R}{\mathbb{R}}
\newcommand{\OK}{\mathcal{O}_K}

% --- Title block ------------------------------------------------------------
\title{%
  A Bounded, Epistemically Labeled Verification Harness\\
  for Phenomenological Quantum-Channel Models\\
  of Curved Spacetime\\[0.5em]
  \large Claim-Boundary Enforcement, Evidence-Class Serialization, and Calibration Traceability for a Single-Qubit Curved-Space Ansatz}
\author{%
  Yun Kwansub%
  \thanks{Independent author. Code distribution:
    \href{https://github.com/Flamehaven-Labs/QSOT-Harness}{github.com/Flamehaven-Labs/QSOT-Harness}.}%
}
\date{2026-06-12 \\ \small QSOT-Harness (reference code commit \texttt{26c0680})}

% --- Document --------------------------------------------------------------
\begin{document}
\maketitle

\begin{abstract}
We present \texttt{QSOT-Harness}, an open-source Python package that implements a phenomenological model mapping classical spacetime curvature metrics (specifically the Frobenius norms of the Riemann and Ricci tensors) to single-qubit quantum decoherence channels. We emphasize that this mapping is a conceptual ansatz without a first-principles physical derivation. The project's main focus is mathematical consistency verification of the model's assumptions and regression testing rather than discovering new physics. 

The central contribution is not a physical result but a \emph{verification architecture} that makes the software's epistemic status machine-readable. Five elements are enforced and serialized in the output: (i) a fixed \emph{claim boundary} that disclaims external physical validity and any new-law claim; (ii) an \emph{evidence class} attached to every reported observation (mathematical invariant, phenomenological model output, policy-derived classification, synthetic-harness output, external-review signal, optional-engine check, or runtime-dependency status); (iii) explicit \emph{surfacing of environment-dependent backends}, recording which audit/governance backend ran and whether it was external or a deterministic mock; (iv) a \emph{calibration manifest} that lists every free or hand-set parameter together with what it governs and where its effect is realized; and (v) a strict \emph{separation of detector self-tests from model-generated signals}. These are exercised by an 8-phase pipeline of 50 checks --- temporal-state axioms, flat and curved baselines, relativistic boosts, a flat-relative Kirkwood-Dirac comparison, a Transfer-Tensor non-Markovianity proxy, a Rust vector-search sidecar, and a bounded compliance review --- run over six spacetime backgrounds as a case study. We document the software's exact limitations, including the arbitrary selection of the sensitivity parameter $\alpha$ and the purely engineering-focused nature of the Rust sidecar. The code is fully reproducible and CI-tested with $94.48\%$ coverage (64 passing tests). The end-to-end reference run reports an overall \texttt{DEGRADED\_PASS} verdict for two explicitly documented, non-failure reasons: the Kirkwood-Dirac optimizer does not converge within its fixed step budget, and the optional governance backend is absent in the baseline environment.

\noindent\textbf{Verification.}
The repository provides a one-command verifier and frozen JSON/Markdown evidence files. The exact execution commands and verification details are provided in \S\ref{sec:repro}.
\end{abstract}

\clearpage
\tableofcontents
\clearpage

% ---------------------------------------------------------------------------
\section{Background and scope}
\label{sec:background}

% 00_theory_intent.tex
% Theoretical Intent and Positioning (QSOT-Harness)

\subsection{Theoretical Intent and Positioning (QSOT-Harness)}
Rather than proposing or claiming the proof of a new fundamental law of quantum gravity, QSOT-Harness introduces a formal \textbf{Temporal-State Axiom Contract Layer} (designated as Phase 0). This layer mathematically evaluates whether temporal quantum evolution can be represented as a constrained multipartite temporal-state structure, rather than treating time strictly as an external classical parameter. 

The primary goal of this implementation is to establish a rigorous \textit{falsification harness} for the candidate Time-as-State hypothesis under explicitly bounded model assumptions. It enforces mathematical constraints such as Kraus-map linearity, CPTP completeness, actual trace preservation, trajectory consistency (conditionability), and density matrix validity. This ensures that any conceptual claim regarding temporal-state representation remains executable, testable, and falsifiable in a software environment.

% Section 1: Background and scope

\subsection{Physical context and controlled phenomenological model scope}
The study of quantum fields propagating on curved spacetimes lies at the boundary of general relativity and quantum mechanics. A central challenge in this domain is understanding how classical gravitational fields induce decoherence in localized quantum systems. 

To explore this, \texttt{QSOT-Harness} implements a controlled phenomenological model that maps the geometry of a classical spacetime background (using the Riemann tensor's Frobenius norm for decoherence modeling, and the Ricci tensor for the admissibility gate checks) to a single-qubit quantum channel. We emphasize that this model is a phenomenological ansatz and does not represent a first-principles physical derivation. The model is conceptually inspired by the established idea that gravitational time dilation and spacetime curvature induce decoherence in localized quantum probes (see for instance~\cite{Pikovski2015}), but the specific mathematical mapping is a software construct designed for demonstration.

\subsection{The transition to V2: software verification}
The previous iteration of this software (V1) combined these physics models with a Streamlit web dashboard and a FastAPI REST server. 

The current release, \texttt{QSOT-Harness}, shifts to a print-and-exit command-line interface (CLI). Rather than acting as an external physical discovery tool, it is designed as a mathematical and software verification harness. The core execution engine is organized into a sequential 8-phase pipeline containing 50 distinct checks. These checks do not validate a law of nature; instead, they verify that the encoded backgrounds, curvature residuals, temporal-state channel maps, and audit gates remain mathematically consistent under the stated QSOT assumptions.

\subsection{Scope and explicit boundaries}
This replication package accompanying \texttt{QSOT-Harness} has a narrow, software-focused scope:
\begin{itemize}[leftmargin=*]
  \item we \emph{do not} validate a new physical law or propose a verified model of quantum gravity;
  \item we \emph{do not} claim that the mapping has a physical derivation from path integrals or field theory;
  \item we \emph{do} implement, in pure NumPy, a mathematical transformation mapping a spacetime background field and its Riemann and Ricci curvatures to Kraus CPTP channels;
  \item we \emph{do} write regression tests verifying that the code consistently computes purity decay under curved metrics and zero decay under flat metrics;
  \item we \emph{do} demonstrate a hybrid Python-Rust software architecture using a vector-search sidecar subprocess as an integration check.
\end{itemize}

The remainder of the paper describes the temporal-state axioms and phenomenological calculations (Phases 0--3) (\S\ref{sec:phases_1_3}), non-classicality checks (\S\ref{sec:phases_4_5}), the Rust turbovec integration (\S\ref{sec:phase_6}), the compliance audit framework (\S\ref{sec:scientific_audit}), explicit limitations (\S\ref{sec:limits}), and reproducibility guidelines (\S\ref{sec:repro}).


% ---------------------------------------------------------------------------
\section{Phases 1--3: Flat Minkowski, curvature noise, and relativistic boosts}
\label{sec:phases_1_3}

% Section 2: Phases 1-3: Flat Minkowski, Curvature Noise, and Relativistic Boosts

This section describes the mathematical formulations and verification steps for the first three phases of the execution engine.

\subsection{Phase 1: Flat Minkowski baselines}
Phase 1 establishes the baseline for a 2-level quantum system (qubit) evolving through flat Minkowski spacetime. The density matrix is initialized in a pure state $\rho_0$. The flat metric is defined by $G = \text{diag}(-1, 1, \dots, 1)$, and the Ricci curvature tensor is identically zero, $R_{\mu\nu} = 0$. 

The CPTP channel is the identity map $\mathcal{E}_{\text{flat}}(\rho) = \rho$. The phase executes five regression tests to ensure that the code correctly computes:
\begin{enumerate}
  \item \tname{flat_rest_purity_is_1}: Verifies that the purity $P(\rho) = \text{Tr}(\rho^2)$ remains exactly $1.0$.
  \item \tname{flat_rest_entropy_is_0}: Verifies that the Von Neumann entropy $S(\rho) = -\text{Tr}(\rho \ln \rho)$ remains exactly $0.0$.
  \item \tname{flat_gate_passes}: Confirms that the background verification gate accepts flat Minkowski as a baseline.
  \item \tname{flat_ctc_is_clear}: Confirms the absence of closed timelike curves.
  \item \tname{flat_boost_purity_is_1}: Confirms that moving the observer through flat Minkowski at relativistic speed does not introduce spurious decoherence.
\end{enumerate}

\subsection{Phase 2: Curvature-induced decoherence and the Riemann-Ricci distinction}
Phase 2 implements the curved-space phenomenological mapping. Rather than mapping the Ricci curvature tensor directly to decoherence (which would physically contradict Ricci-flat solutions like Schwarzschild), the model distinguishes between Ricci curvature (representing energy-momentum sources evaluated for gate checks) and Riemann curvature (representing tidal gravitational forces evaluated for decoherence). The Frobenius norm of the Riemann tensor $\|R^\lambda{}_{\mu\nu\rho}\|_F$ is mapped to a decoherence probability $p$:
\begin{equation}
  p = 1 - e^{-\alpha \|R^\lambda{}_{\mu\nu\rho}\|_F},
  \label{eq:curvature_prob}
\end{equation}
where $\alpha \in (0, 100]$ is a free sensitivity parameter. The probability $p$ is used to build the Kraus operators $\{K_i\}$ of a CPTP channel. For a depolarizing channel, these are:
\begin{equation}
  K_0 = \sqrt{1-p}\,\mathbb{I}, \quad K_1 = \sqrt{\frac{p}{3}}\,X, \quad K_2 = \sqrt{\frac{p}{3}}\,Y, \quad K_3 = \sqrt{\frac{p}{3}}\,Z.
\end{equation}
The density matrix evolves according to:
\begin{equation}
  \rho_{n+1} = \sum_i K_i \rho_n K_i^\dagger = (1-p)\rho_n + \frac{p}{3}\sum_{j=1}^3 \sigma_j \rho_n \sigma_j.
\end{equation}

\subsubsection{Resolution of Ricci-flat vacuum states}
Theoretically, both Schwarzschild and Eguchi-Hanson spacetimes are Ricci-flat vacuum solutions ($R_{\mu\nu} = 0$), which means they satisfy the classical Einstein field equations in vacuum (gate PASS). However, they possess non-zero Riemann curvature tensors representing non-vanishing tidal gravitational fields. 

By mapping the decoherence channel to the Riemann norm ($\|R\| \approx 10^{-3}$) rather than the Ricci norm, the package successfully simulates purity decay and entropy growth under Schwarzschild and Eguchi-Hanson spacetimes, while maintaining their correct gate status as Ricci-flat vacuums. This eliminates the previous logical contradiction where artificial Ricci residuals had to be injected to force decay.

\subsection{Phase 3: Relativistic boosts and time dilation}
Phase 3 incorporates observer kinematics and gravity-induced time dilation. When a qubit is observed by a boosted reference frame moving at velocity $v = \beta c$ relative to the metric source, or is situated in a gravitational potential $g_{00}$, the effective time dilation factor is:
\begin{equation}
  \gamma_{\text{total}} = \gamma_{\text{grav}} \cdot \gamma_{\text{boost}} = \frac{1}{\sqrt{-g_{00}}} \cdot \frac{1}{\sqrt{1 - \beta^2}}.
\end{equation}

The time dilation amplifies the effective decoherence probability $p_{\text{eff}}$ experienced by the quantum state over a unit coordinate step:
\begin{equation}
  p_{\text{eff}} = 1 - (1 - p)^{\gamma_{\text{total}}}.
  \label{eq:boosted_prob}
\end{equation}

The verifier executes checks to validate that the code behaves consistently with eq.~\eqref{eq:boosted_prob}, verifying that purity decays faster under a non-zero boost ($P_{\text{boosted}}(\rho_n) < P_{\text{rest}}(\rho_n)$) and that the math simplifies correctly.


% ---------------------------------------------------------------------------
\section{Phases 4--5: Kirkwood-Dirac optimization and non-Markovianity}
\label{sec:phases_4_5}

% Section 3: Phases 4-5: Kirkwood-Dirac Optimization and Non-Markovianity

This section describes the non-classicality checks and memory-effect measurements.

\subsection{Phase 4: Kirkwood-Dirac basis optimization}
The Kirkwood-Dirac (KD) quasiprobability distribution provides a representation of quantum states in terms of two non-orthogonal measurement bases. For a density matrix $\rho$ and two projectors $P_a = |a\rangle\langle a|$ and $P_b = |b\rangle\langle b|$, the KD quasiprobability is:
\begin{equation}
  Q_{\text{KD}}(a, b) = \text{Tr}(P_b P_a \rho).
  \label{eq:kd_formula}
\end{equation}

When the two bases do not commute, $Q_{\text{KD}}(a, b)$ can take complex values whose real part may be negative. We stress that such negativity is a generic property of single-qubit states and is present even in the flat-spacetime baseline; the harness therefore does \emph{not} treat raw KD negativity as a curvature-specific signal or as a contextuality proof.

For each state $\rho$, the framework parameterizes $|a\rangle$ and $|b\rangle$ by Bloch-sphere angles $(\theta_a, \phi_a)$ and $(\theta_b, \phi_b)$ and runs a bounded basis search:
\begin{equation}
  \min_{\theta_a, \phi_a, \theta_b, \phi_b} \text{Re}\left(\text{Tr}[P_b P_a \rho]\right)
  \label{eq:kd_opt}
\end{equation}
with PyTorch's Adam optimizer for a fixed, hand-set budget of 50 steps. In the reference run this search does \emph{not} converge within the budget, so the corresponding check \tname{kd_optimization_converged} is reported as \texttt{DEGRADED\_PASS} rather than \texttt{PASS}.

Rather than reporting raw negativity, the harness records the \emph{flat-relative} quantity
\begin{equation}
  \Delta_{\text{KD}} = Q^{\text{de Sitter}}_{\text{KD}} - Q^{\text{flat}}_{\text{KD}}
  \label{eq:kd_delta}
\end{equation}
(observation \tname{kd_delta}) as a comparative signal under the implemented model. In the reference run $Q^{\text{flat}}_{\text{KD}} \approx -0.123$ and $Q^{\text{de Sitter}}_{\text{KD}} \approx -0.012$, giving $\Delta_{\text{KD}} \approx +0.11$. This is interpreted only as a relative, non-converged proxy within the model.

Finding negative KD values for a single-qubit state is a standard property of quantum mechanics, not evidence about spacetime. Phase~4 therefore demonstrates only that the optimizer runs over Bloch-sphere variables and that the comparative signal is recorded together with its convergence status; it does \emph{not} validate the spacetime mapping's physical reality, prove that gravity is quantized, or constitute a contextuality witness for any external system.

\subsection{Phase 5: Transfer Tensor Method and accessibility score}
Phase 5 exercises a non-Markovianity detector and computes a unified compliance score. The detector is run in two clearly separated modes, which earlier versions conflated.

\subsubsection{Transfer Tensor Method memory-kernel proxy}
While the true Transfer Tensor Method (TTM)~\cite{CerrilloCao2014} reconstructs the non-Markovian memory kernel $\mathcal{K}(t_n, t_k)$ via full process tomography, our framework implements a simple trace-distance proxy:
\begin{equation}
  \text{nm\_measure} = \sum_{n=1}^N \|\rho_{\text{actual}}(t_n) - \mathcal{E}_n(\rho(t_{n-1}))\|_{\text{tr}},
  \label{eq:ttm_proxy}
\end{equation}
where $\|A\|_{\text{tr}} = \frac{1}{2}\text{Tr}\sqrt{A^\dagger A}$ is the trace distance and $\mathcal{E}_n$ is the configured CPTP channel applied at step $n$ --- i.e.\ the Markovian one-step prediction, not a fitted generator. The proxy measures the accumulated drift of the actual trajectory from the channel-predicted Markovian path.

\subsubsection{Detector self-test versus model trajectory}
We separate two categories explicitly:
\begin{itemize}[leftmargin=*]
  \item \emph{Synthetic self-test} (\tname{memory_kernel}, evidence class \texttt{synthetic\_harness\_output}): an artificial coherence backflow is injected into the trajectory so that $\text{nm\_measure} > 0$ (reference value $\approx 1.4\times10^{-3}$). A pass here validates only that the detector registers backflow; it is \emph{not} evidence that the model exhibits non-Markovian physics.
  \item \emph{Model trajectory} (\tname{memory_kernel_model_trajectory}, evidence class \texttt{phenomenological\_model\_output}): when the configured channel is applied consistently with no injected revival, the phenomenological map is Markovian by construction and the model's own trajectory yields $\text{nm\_measure} = 0$.
\end{itemize}
The headline non-Markovianity ``pass'' therefore refers to the detector self-test on synthetic input; the model itself produces no non-Markovian backflow.

\subsubsection{Accessibility score}
The unified accessibility score integrates the check statuses and penalties into a single metric $A \in [0, 1]$:
\begin{equation}
  A = A_{\text{base}} - \text{Penalty}_{\text{Markov}} - \text{Penalty}_{\text{KD}} - \text{Penalty}_{\text{Risk}},
  \label{eq:accessibility}
\end{equation}
This is a heuristic compliance metric designed to demonstrate automated software scoring, not a fundamental physical parameter.


% ---------------------------------------------------------------------------
\section{Phase 6: Rust turbovec sidecar vector search verification}
\label{sec:phase_6}

% Section 4: Phase 6: Rust turbovec Sidecar Vector Search Verification

This section describes the integration and verification of the Rust sidecar.

\subsection{Engineering purpose of the Rust sidecar}
Phase 6 executes a compiled Rust binary named \tname{rust_verify} (built from the sources in \tname{src/qsot_v2/rust_sidecar}) as a subprocess. The data exchange uses temporary JSON files (\tname{test_vectors.json} and \tname{rust_verify_result.json}).

We clarify that this phase is purely a hybrid software architecture test. It verifies that the Python runner can spawn the Rust executable, exchange JSON data, and that the Rust vector index (using the \tname{turbovec} library) functions within its own quantization error limits. This phase does not contribute to the physical validity of the spacetime curvature-to-channel mapping; it is an integration and performance sanity check for the package's software stack.

\subsection{Automated checks}
The phase executes four regression tests:
\begin{enumerate}
  \item \tname{rust_verify_binary_runs}: Confirms that the subprocess executes successfully and exits with code $0$.
  \item \tname{rust_turbovec_verdict_is_pass}: Verifies that the overall verdict returned by the Rust binary is \texttt{PASS}.
  \item \tname{rust_turbovec_ingested_correct_count}: Confirms that the sidecar ingested the expected number of vectors.
  \item \tname{rust_turbovec_quantization_error_within_bounds}: Verifies that the difference between the exact dot products and the quantized index query results remains within the index limits:
        \begin{equation}
          \max |S_{\text{exact}} - S_{\text{quantized}}| \le 0.15.
          \label{eq:quantization_error}
        \end{equation}
\end{enumerate}


% ---------------------------------------------------------------------------
\section{Phase 7: Scientific compliance audit framework}
\label{sec:scientific_audit}

% Section 5: Phase 7: Scientific Compliance Audit Framework

This section describes the automated auditing of physical reports.

\subsection{Verification of metadata consistency}
Phase 7 reviews the background verification reports generated by the runner against the project's own predefined expected-behavior table, producing a bounded, environment-dependent review signal --- not external ground truth. It is a surface/packaging sanity check on whether the simulated gate outputs (PASS/FAIL) and the calculated residuals (such as Ricci norms and metric-signature flags) match those internally defined expectations.

The review backend (\tname{spar_framework}) is an external, optionally-installed package, so the audit is explicitly environment-dependent: when it is present the harness uses it, and when it is absent a local deterministic mock runs (both implemented behind a single adapter in \tname{scientific_audit.py}). The result records this in \tname{audit_context} --- \tname{backend_mode} is \texttt{external} or \texttt{mock}, alongside raw and unit-normalized scores and the runtime-carrier status --- so a reader can tell which backend produced each review signal. Every per-background review carries the evidence class \texttt{external\_review\_signal}.

\subsection{Expected audit results}
We emphasize that the gate and verdict outputs below are \emph{policy-derived classifications} produced by the implemented admissibility rules, not direct physical observables: a \texttt{PASS}/\texttt{FAIL} gate or an \texttt{ACCEPT}/\texttt{MINOR\_REVISION} verdict encodes rule compliance under the stated model, not a measured property of nature. The audit reviews six backgrounds and issues a verdict (\texttt{ACCEPT}, \texttt{MINOR\_REVISION}, or \texttt{REJECT}) for each. The regression tests check that:
\begin{enumerate}
  \item \tname{scientific_audit_accepts_flat}: Minkowski baseline is accepted (\texttt{gate = "PASS"}, \texttt{verdict = ACCEPT}).
  \item \tname{scientific_audit_accepts_schwarzschild}: Schwarzschild is audited as a Ricci-flat GR solution, yielding \texttt{gate = "PASS"} and \texttt{verdict = ACCEPT} (consistent with zero Ricci curvature, while decoherence is successfully mapped from the non-vanishing Riemann curvature, resolving any logical contradictions).
  \item \tname{scientific_audit_accepts_ads5}: AdS5 is accepted (\texttt{gate = "PASS"}, \texttt{verdict = ACCEPT}).
  \item \tname{scientific_audit_confirms_de_sitter_expected_fail}: de Sitter has positive Ricci curvature and fails the admissibility gate, yielding \texttt{gate = "FAIL"} and \texttt{verdict = MINOR\_REVISION}.
  \item \tname{scientific_audit_confirms_eguchi_hanson_expected_revision}: Eguchi-Hanson fails the gravity-sector check and is audited as \texttt{gate = "FAIL"}, yielding \texttt{verdict = MINOR\_REVISION}.
  \item \tname{scientific_audit_confirms_godel_universe_expected_fail}: Gödel universe contains closed timelike curves (CTCs) and fails the causality gate check, yielding \texttt{gate = "FAIL"} and \texttt{verdict = MINOR\_REVISION}.
\end{enumerate}

These checks are not limited to confirming that a compliance wrapper executes. They verify whether the implemented spacetime backgrounds, curvature residuals, temporal-state channel maps, and audit gates remain mathematically consistent under the stated QSOT assumptions. They do not, however, constitute external physical validation or proof of a new physical law.


% ---------------------------------------------------------------------------
\section{Limits and explicit non-claims}
\label{sec:limits}

% Section 6: Limits and explicit non-claims

To maintain scientific integrity and prevent overclaiming, this section documents the explicit limits and boundaries of the \texttt{QSOT-Harness} verification artifact.

\subsection{Phenomenological approximation, not fundamental QG}
The curvature-to-channel mapping:
$$p = 1 - e^{-\alpha \|R^\lambda{}_{\mu\nu\rho}\|_F}$$
is a phenomenological model. It is \emph{not} a first-principles derivation from loop quantum gravity, string theory, or semiclassical gravity path integrals. The framework simulates the consequences of such a mapping, but does not prove the mapping's physical reality.

\subsection{Sensitivity as a free calibration parameter}
The parameter $\alpha$ (sensitivity) is a free calibration parameter. It is not derived from fundamental constants (such as Newton's constant $G$ or Planck's constant $\hbar$). The value $\alpha = 0.1$ used in the default configuration is chosen for numerical stability during demonstration and testing.

\subsection{Resolution of the Schwarzschild Ricci-flatness}
Theoretically, both Schwarzschild and Eguchi-Hanson spacetimes are Ricci-flat GR solutions, meaning their Ricci tensor is zero, $R_{\mu\nu} = 0$. In earlier versions, this introduced a logical contradiction where artificial Ricci residuals had to be injected to force decay, conflicting with the Phase 7 audit. 

This contradiction has been resolved in the code by separating Ricci curvature (representing energy-momentum sources used for gate admissibility checks) and Riemann curvature (representing tidal gravitational forces used for the decoherence mapping). Schwarzschild and Eguchi-Hanson now correctly have zero Ricci norms (passing the gate checks and beta verifier) but non-zero Riemann norms, inducing consistent physical decoherence.

\subsection{KD negativity as a bounded, flat-relative sanity check}
The Kirkwood-Dirac basis search
$$\min_{\theta_a, \phi_a, \theta_b, \phi_b} \text{Re}\left(\text{Tr}[P_b P_a \rho]\right)$$
is a bounded numerical sanity check, not a proof of non-classicality. Single-qubit KD negativity is generic and appears even in the flat baseline, so raw negativity is not curvature-specific. The harness therefore reports only the flat-relative difference $\Delta_{\text{KD}} = Q^{\text{de Sitter}}_{\text{KD}} - Q^{\text{flat}}_{\text{KD}}$ as a comparative signal under the implemented model, and labels the optimizer's (non-)convergence explicitly. It does \emph{not} prove that spacetime is quantized, does \emph{not} reconstruct a multi-time process matrix, and is \emph{not} a contextuality witness for any external system.

\subsection{TTM memory kernel proxy limits}
The Transfer Tensor Method calculation uses a discrete trace-distance-based proxy:
$$\text{nm\_measure} = \sum_{n=1}^N \|\rho_{\text{actual}}(t_n) - \mathcal{E}(\rho(t_{n-1}))\|_{\text{tr}}$$
to quantify the deviation from Markovian evolution. A true non-Markovian characterization requires full process tomography to build the Choi-Jamio{\l}kowski state of the channel, which is outside the scope of a single-qubit sequential simulation. Moreover, the non-zero $\text{nm\_measure}$ reported by the harness originates from a synthetic injected-backflow \emph{self-test}; the model's own trajectory is Markovian by construction ($\text{nm\_measure} = 0$), so a detector pass is not a claim of non-Markovian physics in the model.

\subsection{Compliance audit boundaries}
The scientific compliance audit is a bounded, environment-dependent review signal that checks whether the simulation outputs match the project's own predefined expected-behavior table --- not external ground truth. Its gate and verdict labels are policy-derived classifications, not physical observables, and it does not perform formal verification of the general-relativity metric tensors or solve the Einstein field equations.

\subsection{Inspiration reference, not a derivation source}
The reference \cite{Pikovski2015} is cited only to acknowledge the established physical idea that gravitational time dilation and spacetime curvature induce decoherence in localized quantum systems. This software is an independent phenomenological construct and does not reproduce any specific formula, theorem, or numerical result from that work; in particular, the Riemann-norm-to-Kraus mapping of eq.~\eqref{eq:curvature_prob} is our own software ansatz and is not derived from it.


% ---------------------------------------------------------------------------
\section{Reproducibility}
\label{sec:repro}

% Section 7: Reproducibility

This section provides the complete execution surface to reproduce all numerical findings and test results reported in this paper.

\subsection{Environment setup}
To isolate dependencies, all commands must be executed within the project's activated Python virtual environment:
\begin{verbatim}
  # Windows PowerShell
  .\venv\Scripts\Activate.ps1
\end{verbatim}
Install the package in editable mode with development and PyTorch dependencies:
\begin{verbatim}
  pip install -e .[dev,torch]
\end{verbatim}

\subsection{Compiling the Rust sidecar}
The Rust sidecar binary must be compiled before executing the tests or the main experiment:
\begin{verbatim}
  cd src/qsot_v2/rust_sidecar
  cargo build --release
  cd ../../..
\end{verbatim}
This produces the compiled executable \tname{rust_verify} at \tname{src/qsot_v2/rust_sidecar/target/release/rust_verify.exe} (on Windows).

\subsection{Executing the test suite}
The test suite consists of 64 unit and integration tests under \tname{tests/}. Run pytest with coverage:
\begin{verbatim}
  python -m pytest --cov=qsot_v2 --cov-report=term-missing
\end{verbatim}
The expected output is:
\begin{verbatim}
  ======================== 64 passed, 1 warning in 12.03s ========================
  TOTAL                                      1432     79    94%
  Required test coverage of 90% reached. Total coverage: 94.48%
\end{verbatim}
These 64 passing tests represent software regression coverage, verifying code execution paths rather than proving physical principles.

\subsection{Executing the end-to-end experiment}
To run the full 8-phase experiment, execute the main entrypoint script:
\begin{verbatim}
  python run_experiment.py --config configs/experiment.yaml
\end{verbatim}
This runs the simulation over all six backgrounds, executes the Rust sidecar, and runs the compliance audit. The expected console summary is:
\begin{verbatim}
  [1/8] Phase 0: Temporal Axioms ... PASS (PASS: 5/5)
  [2/8] Phase 1: Flat Baselines ... PASS (PASS: 5/5)
  [3/8] Phase 2: Curvature Noise ... PASS (PASS: 15/15)
  [4/8] Phase 3: Boosts ... PASS (PASS: 5/5)
  [5/8] Phase 4: Kd Governance ... DEGRADED_PASS (PASS: 3/5)
  [6/8] Phase 5: Ttm Accessibility ... PASS (PASS: 5/5)
  [7/8] Phase 6: Rust Turbovec ... PASS (PASS: 4/4)
  [8/8] Phase 7: Scientific Audit ... PASS (PASS: 6/6)
  ----------------------------------------------------------------------
  ✓ Overall Verdict: DEGRADED_PASS
    Passed: 48 / Failed: 0 / Skipped: 1 / Degraded: 1
\end{verbatim}

The overall \texttt{DEGRADED\_PASS} verdict arises from two explicitly documented, non-failure conditions rather than any error:
\begin{enumerate}
  \item \tname{kd_optimization_converged} is \texttt{DEGRADED\_PASS}: the Kirkwood-Dirac basis search does not converge within its fixed 50-step budget (\S\ref{sec:phases_4_5}), so the comparative signal is reported as a non-converged proxy and the check is degraded rather than passed.
  \item \tname{dqe_covenant_validated} is \texttt{SKIPPED}: the optional package \tname{flame_torch} is not installed in the baseline environment.
\end{enumerate}
Both represent expected software/configuration states (a non-converged numerical search and an absent optional backend), not a physical degradation. The reference run therefore reports \texttt{Passed: 48 / Failed: 0 / Skipped: 1 / Degraded: 1}.

\subsection{Output reports}
The execution generates two files in the output directory (default: \tname{reports/}):
\begin{itemize}[leftmargin=*]
  \item \tname{result.json}: Standardized JSON output containing the exact checks dict, compliance status mapping, and numerical observations.
  \item \tname{report.md}: Human-readable Markdown summary documenting the findings.
\end{itemize}
These files are automatically generated and contain no raw, internal jargon.


% ---------------------------------------------------------------------------
\clearpage
\section*{Acknowledgements}

% Section 8: Acknowledgements

The author is grateful to the Google DeepMind team and the engineering group behind Advanced Agentic Coding for the pair-programming and refactoring support. We also acknowledge the contributors to the NumPy, PyTorch, and Rust open-source communities for maintaining the foundational tools that make numerical physics simulations and reproducible verification pipelines possible.


% ---------------------------------------------------------------------------
\clearpage
\printbibliography

\end{document}
